% Transcription par lecture visuelle directe, sans OCR.
% Statut : transcrit ; relecture indépendante à effectuer.
% Source : PDF original, vue 126, côté droit, page manuscrite 253, après le chapitre dixième.
\chapter{Deuxième suite de la préfecture de M. Franchet. Situation de la Congrégation pendant les années de persécution. Vice préfecture.}

\section*{M. Franchet préfet. Gouvernement de M. Monnier, 2\up{e} assistant}
% Mentions marginales de la page manuscrite 253.
Nous lisons dans l'histoire Romaine qu'après la perte d'une bataille désastreuse, dans laquelle un consul avoit péri, l'armée avoit été presque anéantie, et dont le résultat mettoit la ville de Rome même dans un danger imminent, le Sénat et le peuple Romain adressèrent au consul survivant, resté à la tête des tristes et minces débris de l'armée vaincue, des lettres de félicitation, de ce qu'il n'avoit pas désespéré du salut de la république. La Congrégation presque expirante auroit pu au moment de l'arrestation de son préfet et de son premier assistant, en se voyant ainsi menacée d'une inévitable dissolution, adresser les mêmes actions de grâces à son second assistant, le seul de ses chefs demeuré sur le champ de bataille, car lui aussi ne désespéra pas du salut de la chose publique. Ce second assistant étoit M. Monnier. Ceux d'entre vous, Messieurs, qui ont un peu vieilli dans la Congrégation doivent avoir, tout comme moi, conservé un tendre souvenir pour cet excellent ami, dont une mort prématurée a terminé trop tôt la brillante carrière. Peut-on se rappeler sans éprouver un vif attendrissement cette douceur angélique, cette candeur d'enfant qui excitoient si vivement toutes les
% Source : PDF original, vue 127, côté gauche, page manuscrite 254.
sympathies et donnoient un charme indéfinissable à tous les rapports qu'on pouvoit avoir avec lui ? Aux qualités aimables du cœur, il unissoit à un degré éminent les connoissances d'un esprit très cultivé et la brillante élocution d'un orateur distingué. Aussi étoit-il parvenu bien jeune à occuper un rang élevé dans la magistrature. Mais il ne se bornoit pas à se faire chérir par ses vertus sociales, et estimer par son éloquence d'homme de loi ; -- Dieu avant tout ! Telle étoit sa devise. Sa piété étoit tendre et expansive ; elle avoit marqué sa place dans les rangs de la Congrégation ; jusqu'au moment où il a plu au Seigneur de l'appeler à lui, il nous y a constamment édifiés par ses vertus.

Ce fut donc sur la tête de M. Monnier que l'autorité tout entière se trouva concentrée, à cette époque critique où plus que jamais le secours d'en haut étoit nécessaire à la Congrégation. Comment auroit-elle pu résister à un choc aussi terrible si elle n'avoit été visiblement soutenue par la protection du Tout puissant ?

\section*{Mesures de prudence}
% Mention marginale de la page manuscrite 254 ; reprise aux pages 256 et 257.
M. Monnier se mit de suite en possession du pouvoir et ordonna toutes les mesures de prudence que nécessitoit l'état des choses. Les circonstances étoient déjà fort graves, mais tout pouvoit faire craindre qu'elles ne le devinssent encore davantage. La plupart des Congréganistes ignoroient la véritable cause de l'arrestation de nos chefs ; ils pouvoient supposer assez naturellement que la persécution dirigée d'abord contre eux s'étendroit plus tard à tous ceux qui composoient la Congrégation. Les rapports que M. Franchet avoit eus avec le commissaire général de police et le refus qu'il avoit fait de livrer nos noms étoient des motifs bien suffisants pour le faire présumer.
% Source : PDF original, vue 127, côté droit, page manuscrite 255.
Ceux même qui, dans le secret de ce qui s'étoit passé, connoissoient les véritables causes de ces arrestations, ne pouvoient pas non plus être sans inquiétude, car ils savoient assez combien le gouvernement étoit ombrageux et toujours prêt à envelopper de nouvelles victimes dans ses mesures de rigueur sur le moindre soupçon. Il n'eut donc point été étonnant qu'il fît peser le joug de la tyrannie sur une société qui n'avoit pas voulu se soumettre à ses exigences. Très heureusement il croyoit que cette société n'existoit réellement plus : s'il eut connu la vérité, il ne nous auroit certainement pas épargnés.

Tout nous faisoit donc alors une loi de nous renfermer plus que jamais dans notre système de prudence, et de faire disparoître entièrement tous les indices qui auroient pu déceler que la Congrégation subsistoit toujours. C'est pourquoi le second assistant se hâta de faire donner promptement connoissance de l'arrestation de nos amis à tous les Congréganistes, et fit accompagner cette communication de l'ordre positif de détruire à l'instant même tous les papiers relatifs à la Congrégation qu'ils pourroient avoir chez eux ; toutefois à l'égard de ceux dont la conservation paroîtroit trop essentielle on les autorisoit à les conserver, mais en les cachant de telle manière qu'ils fussent parfaitement à l'abri de toute recherche.

Tout le monde comprenoit alors combien ces mesures étoient nécessaires ; on exécuta donc les ordres reçus avec une exactitude rigoureuse. Il faut même l'avouer, car l'histoire doit être le miroir fidèle de ce qui s'est passé, la terreur avoit tellement gagné les esprits en présence d'événements aussi pénibles, qu'elle dégénéra même en pusillanimité. Les parents des Congréganistes surtout qui craignoient de se voir enlever leurs enfants, par suite des découvertes que la police pourroit faire chez eux,
% Source : PDF original, vue 128, côté gauche, page manuscrite 256.
exigèrent qu'on fît main basse sur tout ce qui pouvoit les compromettre ; aussi malgré l'exception formelle qu'avoit fait le second assistant en faveur des papiers essentiels, règlements et procès verbaux dont la conservation eut été si facile à concilier avec la prudence, tout fut impitoyablement livré aux flammes. C'étoit un véritable holocauste offert en sacrifice à la peur. Cette perte, celle des procès verbaux surtout, a été vivement sentie ; nous étions si heureux d'y retrouver retracés souvent en traits de feu, les émotions si douces que nous avoit fait éprouver dans un meilleur temps la célébration des solennités de la Congrégation ; nous avons, par bonheur, pu retrouver ces souvenirs dans les cœurs des Congréganistes, ils y avoient produit dans le temps une impression trop vive pour pouvoir jamais en être effacés.

Tout en prenant ces mesures de précaution, on eut le soin de prévenir les Congréganistes, que la Congrégation malgré la tempête, ne cesseroit pas d'exister ; que bien au contraire, plus les circonstances paroissoient difficiles, plus nous devions nous efforcer de resserrer autour de Marie notre bonne mère, les nœuds qui nous unissoient, qu'il falloit donc plus que jamais recourir à sa protection et nous confier pleinement à notre auguste patronne, dont l'appui ne nous manqueroit pas. C'est à son ombre tutélaire que nous espérions nous mettre à l'abri et nous cacher pour laisser passer l'orage. Tous les membres furent spécialement invités à prier pour les pauvres captifs.

Pour concilier la sûreté des Congréganistes avec la conservation de la Congrégation, il fut décidé qu'à l'avenir les sections ne s'assembleroient plus que par fractions de trois ou quatre
% Source : PDF original, vue 128, côté droit, page manuscrite 257.
membres au plus, les chefs de sections devoient désigner des vice présidents pour présider ces divisions en leur absence et veiller à ce que leur réunion eut lieu exactement tous les huit jours ; tout cela fut parfaitement exécuté. Ces petites assemblées ont remplacé les assemblées des sections jusqu'à ce que le retour d'un peu plus de sécurité permît de les réunir de nouveau. On ne conservoit ni note, ni liste, ni aucune espèce de papier, tout devoit être classé dans la mémoire. Le petit nombre de membres qui composoient une fraction de section, permettoit de les réunir avec facilité et sans le moindre danger ; on n'étoit point embarrassé pour trouver une salle d'assemblée pour ces réunions ; tout étoit propre à en tenir lieu ; souvent les membres se donnoient rendez-vous sur un quai, sur une place publique, au milieu d'une rue et là tout en se promenant en dépit des nuées d’espions de police dont la ville étoit alors inondée et souvent en leur présence, ils s'entretenoient à voix basse du service du Seigneur et du soulagement du prochain, et recevoient les avis ou les ordres des autorités de la Congrégation. Sans doute les bonnes œuvres auxquelles nos sections pouvoient se livrer étoient alors bien réduites, mais jamais cependant elles n'ont été entièrement interrompues ; ce fut ainsi que les choses se passèrent dans les premiers mois qui suivirent l'arrestation de nos amis. Nous nous étions d'abord flattés qu'ils nous seroient promptement rendus, mais lorsque le second assistant eut acquis la conviction que malgré l'ardeur de nos désirs et les vœux que nous formions pour cette délivrance, il devenoit bien difficile d'assigner un terme prochain à une persécution déjà si prolongée, il ne crut pas pouvoir sans imprudence, laisser reposer plus longtemps le sort de la Congrégation sur l'autorité d'un second assistant demeuré seul pour la gouverner. Le règlement autorisoit en cas d'absence
% Source : PDF original, vue 129, côté gauche, page manuscrite 258.
prolongée du préfet, la nomination d'un vice préfet pour le remplacer. Ce nomination appartenoit au préfet par conséquent dans l'état des choses au second assistant qui jouissoit de tous ses droits. C'étoit pour lui un moyen facile de pourvoir à ce qu'exigeoit la conservation de la société, il s'empressa d'autant mieux d'en user qu'étant lui-même personnellement absorbé par les soins multipliés qu'exigeoit de lui la profession d'avocat qu'il exerçoit alors avec distinction, il espéroit, non seulement se décharger du fardeau si lourd qui lui avoit été imposé d'une manière bien inattendue, mais encore des précédentes fonctions de second assistant ; ainsi par une même décision, il nomma un vice préfet et donna sa démission de la place de second assistant. Le vice préfet ne pouvoit pas s'abuser sur l'importance de ses devoirs ; il falloit qu'il trouvât le moyen de conserver la Congrégation, de la réorganiser et en même temps de préserver ses membres de tous les dangers auxquels pouvoit les exposer leur qualité de Congréganiste. C'est ce double but qui devoit être le mobile de toute son administration ; mais pour qu'il fût possible de le remplir dans les circonstances délicates qui pouvoient survenir, il étoit indispensable qu'il pût prendre alors suivant l'exigence des cas, des déterminations promptes et rigoureuses, sans se trouver arrêté ou gêné dans la marche qu'il croiroit devoir embrasser, par les usages ou par les règlements. L'autorité du Vice Préfet devoit donc être considérée comme affranchie de toutes les règles ; elle pouvoit être assimilée à celle dont la République Romaine aux jours de dangers investissoit ses dictateurs. Pour partager l'immense responsabilité que ce pouvoir absolu faisoit peser sur lui avec des amis capables de guider sa marche dans une voie hérissée d'obstacles de tout genre,
% Source : PDF original, vue 129, côté droit, page manuscrite 259.
le Vice Préfet se hâta de nommer un premier et un second assistants. Il choisit pour remplir ces fonctions deux chefs de section, sans les décharger pour cela de la présidence des sections, à la tête desquelles ils se trouvoient placés, dérogeant ainsi à l’usage établi depuis bien des années, de ne pas comprendre les assistants dans le nombre des membres classés dans les sections. Il importoit dans l'état où se trouvoit alors la Congrégation de concentrer l'autorité sur un petit nombre de têtes, et l'action immédiate qu'exerçoient comme chefs, les assistants sur leurs sections ne pouvoit pas être inutile dans un temps où le salut de la Congrégation exigeoit qu'elle fût gouvernée despôtiquement.

\section*{Nomination d'un vice préfet : M. Benoît Coste}
% Mention marginale de la page manuscrite 258 ; placée après la description de cette nomination.
Le nombre des Congréganistes en activité avoit sensiblement diminué ; d'un côté les réceptions étant suspendues il ne restoit à la Congrégation aucun moyen de remplir les vides que faisoient chaque jour dans ses rangs les changements de domicile ; de l'autre la peur avoit dispersé un certain nombre de membres timides.

Un tout petit peloton étoit demeuré réuni et groupé autour de l'étendard de Marie, du moins ceux qui le composoient étoient bien résolus de ne jamais l'abandonner ; ils n'étoient qu'environ 20 membres, 24 au plus ; mais c'étoit une troupe d'élite ; on auroit pu les comparer aux soldats de Gédéon ; on ne les auroit pas vus s'étendre sur les bords du torrent pour étancher à leur aise la soif dévorante des biens de la terre ; ils s'étoient bornés à prendre un peu d'eau dans le creux de leurs mains, usant en passant des biens de ce monde comme n'en usant pas, suivant le conseil de l'apôtre, par eux-mêmes ils [mot incertain : étoient] foibles\footnote{(Note du transcripteur) Les traces pâles entre « ils » et « étoient » ne permettent pas d'exclure une négation. Voir Q174.}, comme tout ce qui tient à l'humanité, leur fragilité ressembloit à celle d'un vase de terre, mais dans ce vase qui pouvoit être brisé, notre divin maître avoit
% Source : PDF original, vue 130, côté gauche, page manuscrite 260.
placé l'esprit de foi, comme un flambeau lumineux. Cette lumière céleste éclairoit leurs pas chancelants au milieu des dangers de la persécution. Leur confiance en notre bon Sauveur et en sa Sainte mère les ont rendus invincibles. Les Israélites triomphèrent en s'écriant suivant l'ordre qu'ils en avoient reçu : L'épée du Seigneur et de Gédéon. Nos Congréganistes, au moment du combat se fussent écriés : Notre force est en Jésus et en Marie… et ils étoient sûrs de la victoire.

\section*{Administration du vice préfet}
% Mention marginale de la page manuscrite 260 ; placée au début du paragraphe suivant.
Nous pûmes bien nous apercevoir alors que la persécution ranime la foi et réchauffe la dévotion ; c'est une justice que nous aimons à rendre à notre chère Congrégation, jamais peut-être les sentiments de piété et de zèle et le désir ardent de procurer la gloire de Dieu ne l'animèrent davantage qu'à cette pénible époque. Jamais aussi elle ne se montra plus soumise et plus obéissante à tout ce que ses chefs croyoient devoir exiger d'elle. Les Congréganistes n'ignoroient pas que les vagues étoient amoncelées tout autour de notre frêle embarcation ; mais sans s'inquiéter de la tempête, chacun se confioit dans le pilote, se laissoit conduire et se contentoit d'aider de toutes ses forces à la manœuvre. Quant à celui qui placé au gouvernail, étoit chargé de diriger le navire, les yeux fixés sur l'étoile de la mer, il invoquoit son secours ; avec un pareil guide, il étoit sûr de ne pas s'égarer.

\textit{Respice Stellam, voca Mariam.}\footnote{(Note du transcripteur) Littéralement : « Regarde l'étoile, invoque Marie ». Référence à vérifier, Q171.}

C'est parce qu'ils s'appuyoient ainsi sur le secours de notre auguste patronne, que les chefs de la Congrégation osoient user largement des pouvoirs sans bornes qui étoient entre leurs mains. Comme aucune élection ne pouvoit avoir lieu, l'autorité se perpétuoit dans ceux qui en étoient revêtus, seulement le
% Source : PDF original, vue 130, côté droit, page manuscrite 261.
Vice Préfet se bornoit à donner des remplaçants à ceux qu'une absence volontaire ou forcée avoit momentanément éloignés de nous. La discrétion la plus sévère sur tout ce qui concernoit la Congrégation étoit prescrite à tous ses membres, non seulement vis à vis de ceux qui lui étoient étrangers, mais aussi vis à vis de nos frères qui avoient jugé à propos de s'en séparer ; comme tout le monde ils devoient croire la Congrégation anéantie.

On comprend que ce n'est qu'en agissant avec une prudence extrême qu'on pouvoit espérer de conserver une société dont le sort dans les ténèbres dont elle s'enveloppoit étoit bien digne d'envie, puisqu'elle se trouvoit complètement associée aux persécutions dirigées contre le chef de l'Église.

La Congrégation s'étoit donc, ainsi que vous le voyez, enfoncée dans la profondeur des catacombes ; elle ne pouvoit plus comme elle l'avoit fait dans les temps de sa prospérité, se manifester au dehors par son action ; elle ne vivoit plus que d'une vie souterraine, et cependant aux yeux de celui qui voit et qui connoît tout, jamais peut-être elle n'avoit été plus vivace ; les souffrances de nos frères captifs avoient attiré sur elle les bénédictions du ciel, et la grâce toute puissante de notre Sauveur la conservoit et la soutenoit. On peut dire qu'elle ressembloit alors à une plante dont les rigueurs de l'hiver ont flétri les rameaux, mais dont les racines sont allées puiser de nouveaux sucs dans les entrailles de la terre ; au printemps cette plante reparoîtra revêtue d'une nouvelle fraîcheur ; bientôt vous la verrez se recouvrir de fleurs et de fruits ; c'est ainsi que la Congrégation se préparoit dans le silence à devenir un jour avec le secours de celui qui depuis son origine, l'a toujours conduite pas à pas, ce que vous voyez devant vos yeux qu'elle est maintenant.

Il importoit pendant ce temps de retraite, d'éviter avec soin tout
% Source : PDF original, vue 131, côté gauche, page manuscrite 262.
ce qui auroit pu la faire soupçonner de vivre encore. On indiquoit bien toujours comme à l'ordinaire, les communions générales, mais chacun les faisoit dans sa paroisse. Cependant lorsque l'anniversaire du troisième dimanche de juillet vint encore une fois faire battre nos cœurs en réveillant nos vieux souvenirs, nous pûmes sans inconvénient profiter de la fête de St Bonaventure pour nous mêler à la foule des fidèles qu'elle réunissoit dans l'église qui lui est consacrée, et y communier tous ensemble. Il étoit doux pour chaque frère, dans ces temps d'exil et de séparation, de pouvoir prier à côté de son frère ; tous les cœurs confondus en un vœu, toutes les prières s'unissoient pour faire une sainte violence au ciel. Il nous sembloit que dans sa miséricorde infinie notre bon maître, pour nous offrir quelques compensations aux sacrifices qu'il exigeoit de nous dans ces temps de privations, répandoit sur nos âmes avec plus d'abondance encore qu'à l'ordinaire, les grâces ineffables qui découlent du sacrement de son amour.

\section*{Célébration du 3\up{e} dimanche de juillet 1812}
% Mention marginale de la page manuscrite 262 ; placée après le récit de la communion.
Ce fut ainsi que se passa la première année de la persécution. Les circonstances n'étoient point changées lorsque nous parvînmes à la seconde. La captivité de nos amis duroit toujours ; la moindre indiscrétion pouvoit attirer les plus grands malheurs sur la Congrégation et sur ses membres ; il étoit donc toujours tout aussi nécessaire de ne pas s'écarter du système de prudence dans lequel on étoit entré ; cependant vous le savez, c'est un effet naturel de l'habitude d'accoutumer l'homme à tout, et lorsque le danger nous menace sans cesse, on finit par se familiariser avec lui ; on se permettoit donc parfois quelques petites licences ; ainsi lorsque l'époque du troisième dimanche de juillet revint encore en 1812 cela nous fit presque oublier que nous vivions dans un état continuel
% Source : PDF original, vue 131, côté droit, page manuscrite 263.
d'alarmes et nous cherchâmes les moyens de pouvoir savourer ensemble les vieux souvenirs de la Congrégation et de nous livrer en secret à quelques épanchements de cœur. Comme l'année précédente la communion générale eut lieu à St Bonaventure, mais l'après midi la Congrégation ne pouvoit pas sans une notable imprudence essayer de se réunir ; on divisa tous les membres qui la composoient en fractions pour en former trois réunions différentes ; on observa de mélanger les membres des trois sections de manière à donner à ceux que les mêmes œuvres ne rapprochoient pas habituellement, le bonheur de se retrouver ensemble au moins une fois en passant. Cette réunion se prolongea pour chacune des trois fractions pendant toute l'après midi et la soirée du jour de notre anniversaire, mais les exercices en usage pour la célébration du troisième dimanche eurent lieu successivement dans chacune des fractions, afin que le vice-préfet pût présider les trois assemblées partielles. Il se multiplioit ainsi pour servir par sa présence au milieu de tous comme de point de centre à la Congrégation tout entière. Il prit la parole dans chacune des réunions pour rappeler à tous les Congréganistes tout ce que nous devions à la puissante intervention de notre patronne et les marques visibles de protection que la Congrégation avoit reçues dans tous les temps et recevoit encore tous les jours de la divine Providence. Un tout petit goûter vint dans chaque réunion donner l'essor à la douce gaîté qui n'abandonnoit pas les Congréganistes, même dans des temps aussi mauvais. La joie y fut très vive ; s'il n'étoit pas possible de déployer dans la célébration de cette belle fête toutes les pompes dont elle est accompagnée en temps de paix, on n'en goûtoit pas moins avec un plaisir indicible
% Source : PDF original, vue 132, côté gauche, page manuscrite 264.
le bonheur bien senti de pouvoir encore presser sur son cœur un certain nombre de ses frères, et dans chacune des réunions le vice-préfet devenoit le représentant et le dépositaire des deux autres ; il paroissoit partout comme l'expression des mêmes sentiments, pour unir ces trois divisions dans une seule pensée, par laquelle toutes les trois n'en formoient réellement qu'une.

Plus nous avancions et moins il étoit possible d'assigner un terme prochain à une persécution qui avoit placé notre Congrégation dans un état si précaire. Cependant nous ne pouvions pas nous dissimuler qu'en se prolongeant encore, cette proscription devoit rendre bien incertaine la durée de notre société ; il étoit évident qu'en peu d'années, l'absence de toutes les réceptions alloit la faire tomber dans un état d'anéantissement complet. En continuant à suivre le même régime nous arrivions donc à mourir d'inanition de peur de nous exposer à une mort violente ; nous recruter ou ne pas nous recruter, devenoit tout à fait pour nous une question de vie ou de mort ; nous espérâmes qu'amis ou ennemis, personne ne doutoit plus que la Congrégation ne fût parfaitement morte et enterrée ; d'ailleurs Napoléon et sa police étoient alors trop occupés des affaires extérieures pour avoir le temps de penser à nous. Nous crûmes donc qu'en agissant toutefois avec notre prudence accoutumée, le moment étoit arrivé de commencer à nous réveiller de notre assoupissement.

\section*{On fait des réceptions}
% Mention marginale de la page manuscrite 264.
Il fut décidé qu'on alloit s'occuper de réceptions, en prenant les moyens pour que ce recrutement devenu nécessaire, ne pût pas nuire à notre sécurité. Depuis longtemps la section du zèle ne se réunissoit plus ; elle reçut l'ordre de s'assembler et de reprendre ses travaux. On dirigea ses opérations de telle manière qu'elle ne proposa à la Congrégation que des sujets parfaitement dévoués à la cause de la
% Source : PDF original, vue 132, côté droit, page manuscrite 265.
religion et assez au dessus de toute pusillanimité pour qu'on pût compter sur leur courage au jour du danger.

Le chef de la section du zèle réclamoit le concours des informations des membres, soit par lui-même, soit par l'intermédiaire des chefs de section. Lorsqu'un sujet étoit proposé par la section du zèle pour être reçu, on réunissoit les diverses fractions de chaque section en une assemblée générale de la section ; le Vice Préfet, les assistants et les chefs de la section du zèle, se rendant successivement dans ces trois assemblées ; les assistants recueilloient tour à tour les suffrages des membres qui les composoient ; ils avoient soin de prendre note dans chaque section du nombre des boules blanches ou noires, et ce n'étoit que lorsque le vote opéré dans les trois sections avoit complété le scrutin qu'on donnoit à la Congrégation connoissance du résultat.

Le sujet admis étoit présenté à la section du zèle, présidée dans ce cas par le vice-préfet, et qui par une fiction assez naturelle étant composée de membres des trois sections, représentoit assez bien la Congrégation. On ne donnoit point au récipiendaire le livre du règlement puisqu'il avoit été brûlé, mais on lui en faisoit connoître les principales dispositions, et on lui faisoit observer les différences que les circonstances avoient forcé d'apporter entre la lettre du règlement et ce qui se pratiquoit alors. La communion du candidat pour sa réception étoit fixée comme à l'ordinaire à la plus prochaine de nos fêtes. Quelques membres étoient désignés conjointement avec les assistants, qui devoient en être les témoins nécessaires, pour entendre la messe dans la même église que le récipiendaire et l'accompagner à la table Sainte.

Le plus tôt possible après sa réception, le nouveau reçu étoit encore admis dans une assemblée de la section du zèle, présidée par le vice-Préfet. Là il renouveloit la consécration qu'il avoit prononcée
% Source : PDF original, vue 133, côté gauche, page manuscrite 266.
aux pieds des autels et recevoit le baiser de charité. Le Vice préfet le conduisoit ensuite dans les trois sections assemblées extraordinairement afin que tous les Congréganistes pussent connoître, féliciter et embrasser leur nouveau frère ; on terminoit ordinairement cette ronde d'amitié par la section dont le nouveau reçu alloit être appelé à partager les travaux.

Qu'elles étoient touchantes, Messieurs, les réceptions de cette époque de persécutions qui étoit aussi une époque de foi ! Une société battue par la tempête que l'orgueil irrité avoit suscitée contre l'Église de J. C., ouvroit ses rangs à un jeune néophyte que n'effrayoit nullement la fureur menaçante de l'impiété. Animé par le zèle de la maison de Dieu, ce nouveau frère éprouvoit que son ardeur pour le bien augmentoit encore à cause des difficultés qu'il falloit vaincre pour pouvoir l'opérer. Tout cela nous paroissoit avoir quelques rapports avec l'initiation d'un nouveau baptisé, qui dans les premiers siècles de l'Église venoit au milieu de l'assemblée des fidèles, prendre la place laissée vacante par les confesseurs de la foi qui souffroient pour le nom de J. C. dans les prisons.

Ce fut vers ce temps, qu'un de nos amis, dont une mort prématurée devant le Seigneur a couronné la vie toute de foi et de charité, vint prendre place parmi nous. La Congrégation ne peut pas oublier Fleury Richard qui nous a si souvent édifiés par son zèle ardent dont la généreuse activité s'étendoit même au delà du possible.

L'année 1813 étoit commencée. Napoléon avoit assez à faire de se défendre contre l'Europe entière armée contre lui. Nous profitions de l'oubli dans lequel les circonstances politiques
% Source : PDF original, vue 133, côté droit, page manuscrite 267.
contribuoient à nous enfoncer, pour nous émanciper tous les jours davantage. Nous cherchions à nous rapprocher le plus que nous le pouvions de l'exécution littérale du règlement. Lorsque le troisième dimanche de juillet vint à reparoître au lieu de trois fractions nous n'en fîmes plus que deux. Nous songeâmes en même temps à rétablir les élections ; la liste des candidats pour l'élection du Préfet dont il n'avoit pas été question depuis 1809 fut donc publiée.

Cela devenoit tout à fait indispensable dès que cela paroissoit possible, car malgré la bonne volonté de tous les membres de la Congrégation, il ne falloit pas en abuser et ils auroient eu lieu de se plaindre avec justice d'être gouvernés par des autorités dont les pouvoirs n'avoient plus rien de légal puisque les limites que le règlement avoit posées, étoient depuis longtemps dépassées, et que la continuation de l'exercice de leurs fonctions ne pouvoit trouver sa sanction que dans la nécessité.

Le Vice Préfet avoit d'abord eu la pensée de faire faire l'élection du Préfet dans les sections ; mais en cas de ballottage cela pouvoit présenter des inconvénients ; il auroit fallu vraisemblablement y revenir à plusieurs fois avant d'obtenir un résultat. Un autre motif d'ailleurs lui faisoit vivement souhaiter de pouvoir réunir toute la Congrégation, car il désiroit lui faire décider par elle-même une question importante. Si on s'en tenoit à la lettre du règlement M. Franchet ne pouvoit plus être candidat puisque il étoit préfet depuis quatre ans ; mais dans la circonstance toute exceptionnelle où on se trouvoit, le règlement étant déjà suspendu sur tant d'autres points devoit-on s'y tenir ? Le vice-Préfet pensoit qu'à la Congrégation seule appartenoit la solution de
% Source : PDF original, vue 134, côté gauche, page manuscrite 268.
cette question. Après avoir tout bien examiné, on crut sans inconvénient pouvoir rassembler la Congrégation dans une maison de campagne près de la ville, le jour de la fête de l'Assomption. Ainsi ce fut sous les auspices de Marie qu'après une interruption si longue la Congrégation se vit pour la première fois réunie en assemblée générale. La question relative à M. Franchet lui fut soumise, les avis furent unanimes ; la Congrégation déclara que dans les circonstances où nous étions encore elle regardoit l'exécution du règlement comme devant être suspendue, qu'en conséquence elle considéroit comme candidats soit le Préfet soit le Vice-Préfet qui ayant présenté la liste n'avoit pas pu s'y porter. Elle décida en même temps que dans le cas où le Préfet viendroit à être renommé, le vice-Préfet continueroit ses fonctions jusqu'à la délivrance et au retour de M. Franchet. Cette décision prise, on procéda à l'élection. M. Franchet fut réélu à une très forte majorité. Nous ne vîmes pas sans attendrissement la Congrégation vouloir ainsi continuer à être associée à la persécution dont son chef étoit la victime ; c'étoit la démonstration la meilleure de l'excellent esprit dont elle étoit animée ; c'étoit en quelque sorte se grouper autour de son préfet, confesseur de la foi, sous [lecture incertaine : les verroux, les chaînes]\footnote{(Note du transcripteur) Deux formulations sont superposées ou corrigées à cet endroit ; la portée de la rature reste à confirmer. Voir Q179.} et venir pour l'aider à les supporter et manifester l'attachement inébranlable que tous ses amis partageoient avec lui pour le chef de l'Église ; car alors personne n'ignoroit que la captivité de nos deux amis n'avoit pas d'autre cause.

\section*{Sa réélection}
% Mention marginale de la page manuscrite 268 ; placée après le récit de la réélection de M. Franchet.
Vers la fin de 1813, nous jouissions déjà dans la position obscure où nous nous étions placés d'une tranquillité suffisante pour songer à revenir enfin sérieusement à l'exécution complète du règlement ; mais pour cela il falloit consulter ses souvenirs pour le rétablir dans son intégrité ; le vice-Préfet, aidé du concours de
% Source : PDF original, vue 134, côté droit, page manuscrite 269, avant le chapitre douzième.
quelques anciens membres, parvint à refaire notre ancien code presque textuellement. Ce travail étoit terminé et alloit être communiqué à la Congrégation lorsque les événements que la Providence avoit fait marcher à pas de géant rendirent son préfet à la Congrégation et avec lui la copie du règlement que la police impériale avoit pris soin de nous conserver. Nous allons maintenant jeter un coup d'œil rapide sur les graves circonstances qui ont renversé la puissance de Napoléon, donné la paix à l'Église, ramené au milieu de nous nos amis victimes de la tyrannie et fait briller la Congrégation d'un nouveau lustre.
